\chapter{Guix Home: Declarative Dotfiles and Userland}
\label{chap:guix-home}

\epigraph{``There are two kinds of developers: those whose dotfiles are a tangled web of broken symlinks in a forgotten GitHub repo, and those who use Guix Home.''}{--- The Unix Ascetic}

\noindent
So far, we have seen how Guix manages package binaries and how Guix System manages the root operating system. But what about your personal home directory?

For decades, developers have managed their dotfiles (\texttt{.bashrc}, \texttt{.gitconfig}, \texttt{.config/nvim}, etc.) using ad-hoc shell scripts, GNU Stow symlink managers, or bare Git repositories. The result is almost always a brittle mess: symlinks point to deleted targets, local edits conflict with Git merges, and setting up a new laptop takes half a day of manual tinkering.

Enter \textbf{Guix Home}.

\section{The Philosophy of Guix Home}

Guix Home applies the exact same declarative, functional paradigm of Guix System to your unprivileged \texttt{\$HOME} directory:
\begin{itemize}
  \item Your user packages, environment variables, dotfiles, and user daemons are declared in a single Scheme file (\texttt{home-configuration.scm}).
  \item Every update creates a new immutable \textbf{Home Generation}.
  \item If an update breaks your shell or editor config, you can \textbf{roll back} your home environment instantly.
  \item It works on \textbf{any} Linux distribution running GNU Guix, not just Guix System!
\end{itemize}

\section{Writing a \texttt{home-configuration.scm}}

Let us inspect a complete, elegant \texttt{home-configuration.scm}:

\begin{lstlisting}[style=guixscheme]
;; home-configuration.scm - Declarative personal userland
(use-modules (gnu home)
             (gnu home services)
             (gnu home services shells)
             (gnu home services version-control)
             (gnu packages)
             (gnu packages emacs)
             (gnu packages version-control)
             (gnu packages rust-apps)
             (guix gexp))

(home-environment
  ;; Packages to install in the user's home profile
  (packages (list emacs-no-x
                  git
                  ripgrep
                  fd
                  bat))

  ;; Services configuring dotfiles and user daemons
  (services
    (list
      ;; Declarative Bash Configuration
      (service home-bash-service-type
        (home-bash-configuration
          (guix-defaults? #t)
          (aliases '(("ll" . "ls -lha")
                     ("gs" . "git status")
                     ("gshell" . "guix shell --pure")))
          (bashrc (list (plain-file "bashrc-custom"
                          "export EDITOR=emacs\nexport HISTSIZE=50000\n")))))

      ;; Declarative Git Configuration (~/.gitconfig)
      (service home-git-service-type
        (home-git-configuration
          (user-name "Alice Hacker")
          (user-email "alice@example.org")
          (extra-config
            '((init ((defaultBranch . "main")))
              (pull ((rebase . #t)))
              (color ((ui . "auto")))))))

      ;; Deploying raw config files into ~/.config/
      (simple-service 'my-editor-config
        home-xdg-configuration-files-service-type
        `(("alacritty/alacritty.toml"
           ,(local-file "configs/alacritty.toml")))))))
\end{lstlisting}

\section{Applying Your Home Configuration}

To instantiate or update your home environment, you simply run:

\begin{lstlisting}[style=guixcli]
guix home reconfigure home-configuration.scm
\end{lstlisting}

When you run this command:
\begin{enumerate}
  \item Guix validates your Scheme records.
  \item It generates the necessary configuration files and places them immutably in \texttt{/gnu/store}.
  \item It atomically updates symlinks in \texttt{\textasciitilde/.config/} and your home directory.
  \item It creates a new Home Generation in \texttt{\textasciitilde/.guix-home}.
\end{enumerate}

\section{Navigating Home Generations and Rollbacks}

Did you make an experimental edit to your shell configuration that broke your terminal startup?

\begin{lstlisting}[style=guixcli]
# List all historical generations of your home environment
guix home list-generations

# Roll back to the previous home generation in 2 milliseconds
guix home roll-back

# Switch directly to Generation 5
guix home switch-generation 5
\end{lstlisting}

\begin{tearsbox}[The New Laptop Onboarding Experience]
On traditional systems:
\begin{enumerate}
  \item Buy new laptop.
  \item Spend 6 hours installing packages, copying SSH keys, cloning 14 git repos, debugging Python versions, and configuring shell themes.
\end{enumerate}
With Guix Home:
\begin{enumerate}
  \item Buy new laptop. Install GNU Guix.
  \item Clone your \texttt{home-configuration.scm}.
  \item Run \guixcmd{guix home reconfigure home-configuration.scm}.
  \item In 45 seconds, your entire personalized computing environment is bit-for-bit identical to your previous machine.
\end{enumerate}
\end{tearsbox}
